% SOURCE_AUTHORITY: sga5_fr_workpass.tex, lines 13301--13321 (LF-normalized unit).
% SOURCE_SHA256: 791F4EFFC5E02832D5D77ED03518C8156D6F07E4C8238B03545DB93D883FBB28
Deste modo, deduzem-se, em particular, homomorfismos de complexos de $\La_n$-módulos:
\begin{align*}
f'_{H_{n,U',X'}} &= \R\Gamma_{X'}(f'_{\La_{n,U',X'}}) : H_{n,U',X'} \longrightarrow H_{n,U',X'} \quad , \\
f'^\vee_{H_{n,U',X'}} &= (f'_{H_{n,U',X'}})^\vee : H_{n,U',X'}^\vee \longrightarrow H_{n,U',X'}^\vee \quad , \\
g_{H_{n,U',X'}} &= \R\Gamma_{X'}(g_{\La_{n,U',X'}}) : H_{n,U',X'} \longrightarrow H_{n,U',X'} \\
g^\vee_{H_{n,U',X'}} &= (g_{H_{n,U',X'}})^\vee : H_{n,U',X'}^\vee \longrightarrow H_{n,U',X'}^\vee \quad ,
\end{align*}

tais que se verificam as seguintes relações:

\begin{align}
\tag{1.4} g_{H_{n,U',X'}} \circ f'_{H_{n,U',X'}} &= f'_{H_{n,U',X'}} \circ \varphi(g)_{H_{n,U',X'}} \\
\tag{1.5} f'^\vee_{H_{n,U',X'}} \circ g^\vee_{H_{n,U',X'}} &= \varphi(g)^\vee_{H_{n,U',X'}} \circ f'^\vee_{H_{n,U',X'}} \quad .
\end{align}

A relação (1.5) mostra (cf. Exp. XI sobre a teoria geral dos traços, 5.2.1, 7.1) que o traço
\[
\Tr^*_{\La_n[G],\varphi}(f'^\vee_{H_{n,U',X'}}) \in \La_n[G]_{\natural,\varphi} \simeq \mathcal{F}_!(G_{\natural,\varphi}, \La_n) \qquad (*)\footnote{Cf. III B 5.12.}
\]

Nosso primeiro objetivo será dar, sob algumas hipóteses suplementares, uma fórmula explícita para $\Tr^*_{\La_n[G],\varphi}(f'^\vee_{H_{n,U',X'}})$ em termos de invariantes locais associados a $f'$, termos que definiremos mais adiante.
